Diese Dissertation präsentiert die Entwicklung und Performance von Spurrekonstruktionsoftware, die für das High-Luminosity-Upgrade des ATLAS-Experiments am CERN optimiert wurde. Mit dem Beginn des High-Luminosity Large Hadron Collider (HL-LHC), auch als Phase-2 bezeichnet, steht der ATLAS-Detektor vor der neuartigen Herausforderung, Teilchenspuren mit bis zu 200 simultanen Proton-Proton-Kollisionen pro Bunch Crossing zu rekonstruieren. Um den Anforderungen gerecht zu werden, hat sich die ATLAS-Kollaboration entschieden, das Acts-Toolkit (A Common Tracking Software) zu nutzen —-- ein gemeinschaftlich entwickeltes Projekt, das experimentunabhängige Spurrekonstruktionsalgorithmen in modernem C++ bereitstellt.

Die beschriebenen Beiträge erweitern die Funktionalität von Acts erheblich und verbessern ihre Leistungsfähigkeit in verschiedenen Bereichen der Spurrekonstruktion. Dazu zählen unter anderem Verbesserungen bei der Spurpropagation, der numerischen Integration, der Navigationslogik, der Ambiguitätsauflösung sowie der Spurrekonstruktion mit Zeitmessungen. Ein besonderes Augenmerk liegt auf Entwicklungen, die die Recheneffizienz und Robustheit der Spurfindung unter HL-LHC-Bedingungen erhöhen.

Neben der Softwareentwicklung enthält diese Dissertation umfassende Leistungsstudien. Die Performance der Spurrekonstruktion wird sowohl am OpenDataDetector (ODD), einem Prototypen eines (HL-)LHC-ähnlichen Silizium-Innerspur-Detektors, als auch am Inner Tracker (ITk) des ATLAS Phase-2 gemessen. Die Ergebnisse zeigen die Wirksamkeit der neuen Entwicklungen in hochkomplexen Kollisionsumgebungen und präsentieren die Effizienz, Auflösung und Rechenleistung der Spurfindung.

Diese Arbeit legt eine robuste Grundlage für die Integration von Acts in den Rekonstruktionsprozess von ATLAS Phase-2 und leistet einen wichtigen Beitrag zu einem wachsenden Softwaretoolkit für zukünftige Experimente der Hochenergiephysik.
